\chapter{gamma 与 3x ledger：目标仓位如何变成真实仓位}

\section{本章要解决的困惑}

entry 触发以后，策略还没有真正下满仓。它要先算目标仓位，再经过 capacity ledger。
当前控制并不复杂：

\[
target\_exposure_t = weak\_overlay\_weight_t \cdot \gamma_{cell_t}.
\]

然后进入 3x 杠杆账本：

\[
actual\_exposure_t
= \min(target\_exposure_t,\ 3-open\_exposure_t).
\]

这就是主体。

\section{base weight 从哪里来}

\code{BASE_WEIGHTS} 是 membership set 到基础权重的表。它不是模型估计器，而是手写映射。

例如：

\begin{lstlisting}
tfi_follow_flat+tfi_long_flat -> 0.5
tfi_follow_flat+tfi_short_flat -> 0.5
tfi_follow_flat+tfi_short_flat+tfi_short_stale25+tfi_event_active -> 2.0
\end{lstlisting}

设当前 membership set 对应：

\[
base\_weight=0.5.
\]

如果 frames 特别高，超过 overlay threshold，当前代码加一个弱 overlay：

\[
weak\_overlay\_weight
= base\_weight + overlay\_boost.
\]

其中：

\[
overlay\_boost =
\begin{cases}
0.5, & frames\_since\_mid\_change \ge 59.8,\\
0, & \text{otherwise}.
\end{cases}
\]

\section{gamma 从 cell 来}

当前默认 \code{core_q70} preset 是：

\begin{longtable}{p{0.30\textwidth}p{0.20\textwidth}p{0.40\textwidth}}
\toprule
cell & gamma & 解释 \\
\midrule
\code{00_none} & 1.25 & baseline core cell \\
\code{10_r5_only} & 0.75 & R5 好但 frames 不高 \\
\code{01_frames_only} & 0.0 & frames 高但 R5 不好，默认不用 \\
\code{11_r5_frames} & 0.75 & R5 好且 frames 高 \\
\bottomrule
\end{longtable}

注意，这组 gamma 不是深度学习，也不是连续控制函数。它是一张小表。

\section{完整手算例子}

假设当前 entry 信号：

\begin{lstlisting}
membership_set = tfi_follow_flat+tfi_long_flat
base_weight = 0.5
frames_since_mid_change = 65
cell = 11_r5_frames
gamma[11_r5_frames] = 0.75
open_exposure_before = 2.80
leverage_cap = 3.0
\end{lstlisting}

因为 frames 超过 overlay threshold：

\[
weak\_overlay\_weight=0.5+0.5=1.0.
\]

目标仓位：

\[
target\_exposure=1.0 \times 0.75=0.75.
\]

但账本已经开了 2.80，剩余 capacity：

\[
free=3.0-2.80=0.20.
\]

真实仓位：

\[
actual\_exposure=\min(0.75,0.20)=0.20.
\]

被剪掉：

\[
clipped\_exposure=0.75-0.20=0.55.
\]

这就是 \code{entries.parquet} 里 \code{requested_exposure}、\code{actual_exposure}、
\code{clipped_exposure} 的来源。

\section{core\_idle01 profile 的区别}

\code{execution_runtime.py} 里有两个 allocator：

\begin{lstlisting}
CapacityLedger
CoreIdle01CapacityAllocator
\end{lstlisting}

普通 \code{CapacityLedger} 只有一条 FIFO capacity：

\[
actual = \min(requested,\max(0,cap-open)).
\]

\code{CoreIdle01CapacityAllocator} 多了一个特殊处理：\code{01_frames_only} 可以作为 idle sleeve，只用剩余空闲 capacity，不挤占 core cells。它的工程含义是：

\begin{quote}
让某些弱结构只吃闲置杠杆，不要影响主结构。
\end{quote}

但这仍然是账本规则，不是预测模型。

\section{常见误解}

\begin{itemize}[leftmargin=2em]
  \item \term{误解一：gamma 是收益预测。}不是。gamma 是 cell 到仓位倍率的表。
  \item \term{误解二：target exposure 就是真实下单。}不是。真实进入结果还要经过 admission 和 capacity。
  \item \term{误解三：3x 表示每笔都 3 倍。}不是。3x 是总 open exposure 上限。
\end{itemize}

\checkpoint{读完本章，你应该能从 membership、cell、gamma、open exposure 手算 actual exposure。}

